\chapter{Method}
\label{chap:method}
\todo{Add a method overview figure} This chapter presents the methodological contribution of the thesis: a
formal framing of signer appearance as an intervenable nuisance factor
(\S\ref{sec:method-notation}), a taxonomy of appearance counterfactuals
together with two synthesis routes that realise them
(\S\ref{sec:method-diffusion}), and training objectives that consume the
counterfactuals either as additional data or as explicit positives
(\S\ref{sec:method-objectives}). The concrete instantiation --- datasets,
backbones, training configurations, and evaluation protocols --- is
collected in Chapter~\ref{chap:setup}.

\section{Problem formulation}
\label{sec:method-notation}
% Formal notation block (2026-08-03): set/space definitions + typed maps in
% the register of Nalmpantis §4.1-4.2 (run-in headers, one object per
% paragraph, numbered display equations), grounded in this thesis's own task.
This section fixes the notation used throughout the chapter: the data we
train on, the encoder under study, the appearance-counterfactual operator
that both synthesis routes instantiate, and the retrieval task on which
every model is evaluated. The losses that bind these objects together are
formalised in \S\ref{sec:method-losses} and the evaluation protocols in
\S\ref{sec:eval-protocols}; here we define only the objects they operate
on.

\paragraph{Videos and labels.}
Let $\mathcal{X}$ denote the space of signing videos. A training corpus is
a finite set of annotated clips
\begin{equation}
\label{eq:dataset}
\mathcal{D} \;=\; \bigl\{(x_i,\, y_i,\, s_i)\bigr\}_{i=1}^{n}
\;\subset\; \mathcal{X} \times \mathcal{Y} \times \mathcal{S},
\end{equation}
where $y_i \in \mathcal{Y}$ is the linguistic label of the video $x_i$ and
$s_i \in \mathcal{S}$ identifies its signer; when convenient we write
$y(x)$ and $s(x)$ for the label and signer of a clip. The label space
$\mathcal{Y}$ differs between our two supervision regimes. On the isolated,
gloss-labelled benchmarks each video clip shows a single sign and
$\mathcal{Y} = \mathcal{G}$, a gloss vocabulary of $C = |\mathcal{G}|$
classes shared across signers. The continuous corpus supplies no gloss
vocabulary: full sentences are the unit of recording, and $\mathcal{Y}$ is
the set of recorded sentence identities. Throughout, the signer is a
nuisance variable: labels in $\mathcal{Y}$ are defined by what is signed,
never by who signs it, and the retrieval tasks of this thesis ask the
representation to match clips across $\mathcal{S}$ at fixed
$\mathcal{Y}$-label.

\paragraph{Encoder.}
Every training setting, whatever its objective, produces a single video
encoder
\begin{equation}
\label{eq:encoder}
f_\theta : \mathcal{X} \to \mathbb{R}^{d},
\qquad z \;=\; f_\theta(x),
\end{equation}
mapping a video clip to one embedding vector, where $\theta$ collects the
trainable parameters --- a video--text head over frozen backbone features,
the fine-tuned backbone itself, or the view-contrastive encoder; the
instantiations, and the dimension $d$ each uses, are given in
\S\ref{sec:setup-training}. Whether $z$ is $\ell_2$-normalised is a
property of the objective and is stated with each loss. After training,
$\theta$ is frozen and every analysis in this thesis --- retrieval,
probing, feature geometry --- operates on the embeddings $z$ alone.

\paragraph{Appearance counterfactual operator.}
The interventions of this thesis are appearance edits: maps
\begin{equation}
\label{eq:cf-operator}
T_a : \mathcal{X} \to \mathcal{X}, \qquad a \in \mathcal{A},
\end{equation}
where $\mathcal{A}$ indexes the available edit types and $T_a(x)$ shows
the signing of $x$ under a changed appearance. The defining intent is that
$T_a$ moves a video clip across the nuisance factor while leaving the linguistic
content intact,
\begin{equation}
\label{eq:motion-preservation}
y\bigl(T_a(x)\bigr) \;=\; y(x) \qquad \text{for all } a \in \mathcal{A},
\end{equation}
while the appearance --- and, for identity-changing edits, the apparent
signer --- differs. Whether Eq.~\eqref{eq:motion-preservation} actually
holds is a property of the synthesis route, not an assumption we grant
ourselves: the mocap-driven avatar renders satisfy it by construction,
because every render is driven by the same captured motion, whereas the
frame-wise diffusion edits satisfy it only insofar as the editor leaves the
motion untouched, an empirical matter we return to when interpreting
results. For the diffusion route, $\mathcal{A}$ is
the four-type edit taxonomy $\{\text{shirt},\, \text{glasses},\,
\text{skin tone},\, \text{signer swap}\}$; for the avatar route,
$\mathcal{A}$ ranges over avatar--viewpoint pairs
(\S\ref{sec:setup-generation}).

\paragraph{Positive sets.}
For a training video clip $x$ we write
\begin{equation}
\label{eq:cf-set}
\mathcal{T}(x) \;=\; \{x\} \,\cup\,
\bigl\{\, T_a(x) \;:\; a \in \mathcal{A}(x) \,\bigr\}
\end{equation}
for its counterfactual set, where $\mathcal{A}(x) \subseteq \mathcal{A}$
are the edits realised for $x$ (not every edit is generated for every
clip). The contrastive objectives of \S\ref{sec:method-losses} are built
on positive sets over a batch: for an anchor $u$ in a batch indexed by
$\mathcal{B}$,
\begin{equation}
\label{eq:positive-set}
P(u) \;=\; \bigl\{\, p \in \mathcal{B} \setminus \{u\}
\;:\; p \sim u \,\bigr\},
\end{equation}
where $\sim$ is the positive-defining equivalence of the run: membership
of one counterfactual set $\mathcal{T}(x)$ (an original and its edits;
counterfactual-contrastive), a shared gloss in $\mathcal{G}$ (gloss-level
contrast across signers), or a shared recording identity (the view-contrastive setting,
where all $K$ views of one recording --- real cameras and synthetic
renders alike --- are positives). Everything in the batch outside $P(u)$
acts as a negative.

\paragraph{Cross-signer retrieval.}
The task is retrieval with the frozen encoder. An evaluation protocol
supplies a query set $\mathcal{Q}$ and a reference gallery $\mathcal{R}$,
each element carrying a label in $\mathcal{Y}$; for a query $q$ the
gallery is ranked by decreasing cosine similarity
$\cos\bigl(f_\theta(q),\, f_\theta(r)\bigr)$, and retrieval at rank $k$
succeeds if a correctly labelled reference appears among the $k$
highest-ranked items:
\begin{equation}
\label{eq:recall-at-k}
\mathrm{R}@k \;=\; \frac{1}{|\mathcal{Q}|} \sum_{q \in \mathcal{Q}}
\mathbf{1}\Bigl[\, y(q) \in \bigl\{\, y(r) : r \in \mathcal{R}_k(q)
\bigr\} \Bigr],
\end{equation}
where $\mathcal{R}_k(q) \subseteq \mathcal{R}$ are the $k$ gallery items
most similar to $q$. The protocols of \S\ref{sec:eval-protocols}
instantiate $\mathcal{Q}$ and $\mathcal{R}$ per corpus --- including a
video-to-text variant in which the gallery holds text prototypes rather
than videos --- and share one structural property: evaluation signers are
unseen during training, and in the continuous-corpus protocol query and
gallery are
additionally produced by disjoint signers, making the retrieval itself
cross-signer.

\section{Synthetic appearance counterfactuals}
\label{sec:method-diffusion}


\subsection{Diffusion-based appearance editing}
We generate appearance counterfactuals for the isolated-benchmark corpus
by editing real training videos frame-by-frame with an open image
diffusion model; the model, compute, and the statistics of the edited
subset are given in \S\ref{sec:setup-generation}. Rather than
augmenting a uniform sample, targets are chosen worst-first
at the class level: every gloss is ranked by its number of distinct training
signers (ascending), with ties broken by a difficulty proxy --- per-class
nearest-centroid gloss accuracy computed on the from-scratch contrastive
baseline's test embeddings, a clustering measure that degrades precisely when
embeddings organise by signer rather than by sign --- and entire classes are
taken in rank order until the budget fills.
\todome{REWRITE PROPOSAL (03.08): the opener above binds the diffusion
route to ``the isolated-benchmark corpus'', but Setup and Results apply
Flux to NGT too (stride-4 NGT edits in \S\ref{sec:setup-generation}; the
single-view ladder arms E4/E6/E7/E10 are Flux-on-NGT) --- the binding is
false and it also makes the route-vs-corpus framing in the abstract too
clean. Proposed replacement for the first sentence: ``The first route
edits real training videos frame-by-frame with an open image diffusion
model; it requires nothing beyond the RGB video itself, so it applies to
any corpus. The model, compute, and edited-subset statistics are given in
\S\ref{sec:setup-generation}.'' Let Setup alone say which route ran
where. Same issue in the avatar subsection below.}
We propose a taxonomy of four independent appearance edit types
(Fig.~\ref{fig:aug_examples}); for each
selected source video we generate up to four edits, applied frame-by-frame:
\emph{shirt} (a clothing swap, the lowest-risk edit --- the video backbone's own training
augmentations already include colour jitter, so this edit's marginal value
over standard colour augmentation is itself an open question), \emph{glasses}
(addition/removal, with some risk of altering facial expression near the
eyes), \emph{skin tone} (the highest-risk edit for hallucination, e.g.\
occasional extra or malformed hand shapes), and \emph{signer swap} (a full
identity change, carrying the highest potential benefit for a
signer-independent task alongside the highest hallucination risk).

\subsection{Motion-preserving avatar counterfactuals}
For the continuous corpus we propose a second, motion-capture-driven
route: the captured 3D motion of a real signer is retargeted onto
synthetic avatars and re-rendered under multiple camera viewpoints
(\S\ref{sec:setup-generation}). Because this route starts from
ground-truth motion rather than a 2D image edit, the signing motion is
preserved by construction, and the route can vary camera viewpoint as
well as appearance --- an axis a 2D edit cannot reach from a single
input view.
\todome{REWRITE PROPOSAL (03.08), same corpus-binding issue as above:
``For the continuous corpus'' ties the avatar route to NGT by
definition. Propose opening with ``The second route is
motion-capture-driven: ...'' (rest verbatim) and closing with ``Unlike
the diffusion route it requires motion-capture recordings, so it applies
only where capture is controlled.''}

\section{Counterfactual-aware training objectives}
\label{sec:method-objectives}
We train encoders with these counterfactuals in three settings ---
two on an isolated, gloss-labelled benchmark and one on a continuous
corpus without a gloss vocabulary; their concrete configurations
(datasets, backbones, dimensions) are given in
\S\ref{sec:setup-training}. In every setting the counterfactuals enter
the objective either as additional labelled data or as explicit
positives matched to their source video; the loss formulations below
cover both mechanisms.

\subsection{Loss formulations}
\label{sec:method-losses}
Three distinct model components are trained in this thesis, and each has
its own objective; Table~\ref{tab:objective_map} maps components to
objectives before the formulations are given. The \emph{video--text
head} trains only its own parameters on top of frozen backbone features
--- the backbone receives no gradient. The \emph{backbone fine-tuning}
runs do the opposite: they update the video backbone itself, with a
disposable linear gloss head on top. The \emph{view-contrastive encoder}
of the continuous corpus again trains only its own parameters over
frozen backbone features, and has no text tower at all.

\begin{table*}[t]
\centering
\footnotesize
\setlength{\tabcolsep}{4pt}
\begin{tabular}{llll}
\toprule
\textbf{Component} & \textbf{What is trained} & \textbf{Objective} &
\textbf{Counterfactuals enter as} \\
\midrule
\makecell[tl]{Video--text head\\\scriptsize (SignCLIP, both towers,\\\scriptsize on frozen Logos/I3D/pose feats)} & \makecell[tl]{head only;\\input features frozen} &
\makecell[tl]{video--text InfoNCE,\\Eq.~\eqref{eq:nce}} &
\makecell[tl]{extra $(v,t)$ pairs (aug-as-data), or\\pair-SupCon
positives, Eq.~\eqref{eq:supcon}} \\
\addlinespace
\makecell[tl]{Backbone fine-tuning\\\scriptsize (Logos MViTv2-S)} & \makecell[tl]{the backbone itself\\($+$ linear
gloss head)} & \makecell[tl]{master loss, Eq.~\eqref{eq:master}:\\CE $+$
one invariance option} & \makecell[tl]{extra labelled clips and/or\\
invariance-term positives} \\
\addlinespace
\makecell[tl]{View-contrastive encoder\\\scriptsize (SignCLIP video tower only,\\\scriptsize on frozen Logos feats)} & \makecell[tl]{encoder only;\\backbone frozen}
& \makecell[tl]{SupCon over views, Eq.~\eqref{eq:supcon};\\video-only
InfoNCE at $K{=}1$} & \makecell[tl]{synthetic views as\\additional
positives} \\
\bottomrule
\end{tabular}
\caption{Which objective trains which component. The SupCon form of
Eq.~\eqref{eq:supcon} is shared by two components; what differs is the
positive-defining equivalence of Eq.~\eqref{eq:positive-set} ---
counterfactual instance for the head and backbone runs,
recording identity for the view-contrastive encoder.}
\label{tab:objective_map}
\end{table*}

\paragraph{Video--text head objective.}
The video--text head is trained with the symmetric InfoNCE
loss~\cite{oordRepresentationLearningContrastive2018}
we take from VideoCLIP~\cite{xuVideoCLIPContrastivePretraining2021}, over frozen
768-d backbone features; only the head's parameters receive gradients.
For a batch of paired video and text
embeddings $(v_i, t_i)$,
\begin{equation}
\label{eq:nce}
\begin{split}
\mathcal{L}_{\mathrm{NCE}} \;=\;
-\frac{1}{N}\sum_{i=1}^{N}\Bigl[&
\log \frac{\exp(v_i \cdot t_i)}{\sum_{j}\exp(v_i \cdot t_j)}\\
&+ \log \frac{\exp(v_i \cdot t_i)}{\sum_{j}\exp(v_j \cdot t_i)}
\Bigr],
\end{split}
\end{equation}
where the text side is the gloss rendered as a caption; following
VideoCLIP, the embeddings are not $\ell_2$-normalised and no temperature is
applied. Synthetic counterfactuals enter in one of two ways: as additional
$(v, t)$ pairs carrying their source video's gloss (aug-as-data), or --- in
the pair-SupCon configurations --- by replacing the video--text loss with
the supervised-contrastive form of Eq.~\eqref{eq:supcon} below, computed
over video embeddings alone with an original and its edits sharing one
instance label (counterfactual-contrastive). 

\paragraph{Backbone fine-tuning objectives.}
All fine-tuning objectives operate on the same representation: the
encoder embedding $z = f_\theta(x)$ of \S\ref{sec:method-notation},
instantiated as the 768-d CLS token the video backbone produces for a
32-frame clip $x$, which is also exactly the feature later re-extracted for
evaluation. A linear head $h$ maps $z$ to logits over the $C=|\mathcal{G}|$ glosses
of the training vocabulary. Every run minimises a total loss of the form
\begin{equation}
\label{eq:master}
\mathcal{L} \;=\; \mathcal{L}_{\mathrm{CE}}
\;+\; \lambda_{\mathrm{cons}}\, \mathcal{L}_{\mathrm{cons}}
\;+\; \lambda_{\mathrm{sup}}\, \mathcal{L}_{\mathrm{sup}}^{\mathrm{gloss}},
\end{equation}
where $\mathcal{L}_{\mathrm{cons}}$ stands for the run's
appearance-invariance term (the cosine pull below, the CLS-level contrast
that replaces it, or the pair of SDDA terms with their own weights),
$\mathcal{L}_{\mathrm{sup}}^{\mathrm{gloss}}$ for the optional gloss-level
contrastive auxiliary, and the runs differ only in which weights are
non-zero and in what the invariance term treats as positives. The four
paragraphs that follow define these ingredients; they belong to the
backbone fine-tuning runs only, not to the video--text head.

\paragraph{Cross-entropy (all backbone fine-tuning runs).}
Standard softmax cross-entropy $\mathcal{L}_{\mathrm{CE}} =
-\log p_{h(z)}(y)$ against the gloss label $y$. In paired runs the
augmented copies are classified too, carrying their source video's label
(augmentation-as-data), so CE is averaged over the $B$ originals and $M$
augmentations jointly.

\paragraph{Appearance-consistency pull (invariance option 1).}
For each original $z_i$ and each of its appearance-augmented copies
$z_i^{a} = f_\theta(T_a(x_i))$, $a \in \mathcal{A}(x_i)$, present in the batch, a
positive-only cosine pull with no negatives:
\begin{equation}
\mathcal{L}_{\mathrm{cons}} \;=\;
\frac{1}{|\mathcal{P}|}\sum_{(i,a)\in\mathcal{P}}
\bigl(1 - \cos(z_i,\, z_i^{a})\bigr),
\end{equation}
weighted by $\lambda_{\mathrm{cons}}=0.5$. This directly asks the encoder
to make $z$ invariant to the edited appearance attribute.

\paragraph{CLS-token supervised contrast (invariance option 2).}
A stronger variant that replaces the pull by adapting the SupCon
objective~\cite{khosla2020supervised}: we keep its loss form over
$\ell_2$-normalised tokens but redefine the positives so that an
original and all its augmentations --- the counterfactual set
$\mathcal{T}(x)$ --- share one \emph{instance} label (the positive set
$P(u)$ of Eq.~\eqref{eq:positive-set}), and every other clip in the
batch is a negative,
\begin{equation}
\label{eq:supcon}
\mathcal{L}_{\mathrm{sup}} \;=\;
\frac{1}{N}\sum_{u} \frac{-1}{|P(u)|} \sum_{p \in P(u)}
\log \frac{\exp(z_u \cdot z_p / \tau)}
{\sum_{v \neq u} \exp(z_u \cdot z_v / \tau)},
\end{equation}
with temperature $\tau = 0.07$, where $u$ and $v$ range over the $N$ clips
in the batch (anchors with an empty positive set are excluded from the
average). Here this term takes the pull's place in the total loss, with the
same weight $\lambda_{\mathrm{cons}} = 0.5$. Unlike the consistency pull, this also
repels different videos, preventing representation collapse.

\paragraph{Gloss-level supervised contrast (optional auxiliary).}
The same objective as Eq.~\eqref{eq:supcon}, but positives are defined by
the gloss label $y$ instead of the source instance: all same-gloss clips in
the batch are pulled together across signers. Off by default
($\lambda_{\mathrm{sup}}=0$; the run that enables it uses
$\lambda_{\mathrm{sup}}=1$).

\paragraph{SDDA-style invariance (invariance option 3, adversarial).}
Following Fu et al.~\cite{fuSignerDiversitydrivenData2024b}, this variant
replaces the cosine pull with two terms: a KL-consistency loss between the
gloss posteriors of an original and its augmented copy,
$\mathrm{KL}\bigl(p_{h(z_i)} \,\|\, p_{h(z_i^a)}\bigr)$ (weight $1$), which
enforces prediction-level rather than embedding-level agreement, and an
adversarial original-vs-augmented discriminator trained on $z$ through a
gradient-reversal layer (weight $10^{-4}$), so that the encoder is
penalised for representations from which the discriminator can still tell
originals from appearance-edited copies.

\paragraph{View-contrastive objective.}
Architecturally, the view-contrastive encoder is the \emph{video tower}
of the same contrastive model used for the head experiments
(Table~\ref{tab:objective_map}), consuming frozen backbone features
through the identical input projection --- but trained from the standard
language-model initialisation rather than from any text-aligned
checkpoint, and with the text tower never called: the continuous corpus
has no vocabulary to align against. It is trained
with the same SupCon form as
Eq.~\eqref{eq:supcon}, with positives defined by the sign-recording
identity (the continuous-corpus label space $\mathcal{Y}$ of
\S\ref{sec:method-notation}): all $K$ views of a sign (real cameras, avatar renders, or
diffusion edits, depending on the arm) share one label, and views of
other signs are negatives.
At $K{=}1$ the positive set is empty and
Eq.~\eqref{eq:supcon} yields no gradient, so the single-view baseline arm
instead trains with a video-only InfoNCE: each clip in the batch is treated
as its own class, and a cross-entropy is taken over the pairwise cosine
similarities of the $\ell_2$-normalised embeddings, scaled by the same
$\tau = 0.07$.
